%% toc: P15 164 164.3,164.4
\section{Proofs for Popper's probability and quantum chapters}
\label{sec:annex-p15-w02b-p15-027}

Notation is that of \cref{def:p15-w02b-p15-027-finite-frequency,def:p15-w02b-p15-027-cyclic-freedom,def:p15-w02b-p15-027-absolute-freedom}.
Throughout, \(0^0=1\).

\begin{annexitem}{annex:p15-w02b-p15-027-finite-calculus}{thm:p15-w02b-p15-027-finite-calculus}{The finite frequency calculus}
\begin{proof}
Put \(n=N(\alpha)>0\). All classes below are subclasses of the finite
class \(\alpha\), so their numbers are finite, and product classes obey
\((\alpha.\beta).\gamma=\alpha.(\beta.\gamma)\).

Item~1. Every element of \(\alpha\) lies in exactly one of
\(\alpha.\beta\) and \(\alpha.\bar\beta\), so
\(N(\alpha.\beta)+N(\alpha.\bar\beta)=n\). Dividing by \(n\) gives the claim.

Item~2. If \(N(\alpha.\beta)>0\), then
\[
 {}_{\alpha}F''(\beta)\cdot{}_{\alpha.\beta}F''(\gamma)
 =\frac{N(\alpha.\beta)}{n}\cdot\frac{N(\alpha.\beta.\gamma)}{N(\alpha.\beta)}
 =\frac{N(\alpha.\beta.\gamma)}{n}={}_{\alpha}F''(\beta.\gamma).
\]
The second form is the same computation with \(\gamma\) in place of
\(\beta\), using \(N(\alpha.\gamma)>0\) and \(\gamma.\beta=\beta.\gamma\).

Item~3. Both forms of item~2 apply. Suppose \(\gamma\) is independent of
\(\beta\) within \(\alpha\), that is
\({}_{\alpha.\beta}F''(\gamma)={}_{\alpha}F''(\gamma)\). Substituting in the
first form of item~2 gives the special multiplication theorem.
Conversely, from the special multiplication theorem and the first form,
\({}_{\alpha}F''(\beta)\cdot{}_{\alpha.\beta}F''(\gamma)
={}_{\alpha}F''(\beta)\cdot{}_{\alpha}F''(\gamma)\); since
\({}_{\alpha}F''(\beta)=N(\alpha.\beta)/n>0\), cancelling gives
independence of \(\gamma\) from \(\beta\). The same two steps with the
second form of item~2, cancelling \({}_{\alpha}F''(\gamma)>0\), show that
independence of \(\beta\) from \(\gamma\) is also equivalent to the special
multiplication theorem. Hence all three conditions are equivalent.

Item~4. Put \(X=\alpha.\beta\) and \(Y=\alpha.\gamma\). Then
\(\alpha.(\beta+\gamma)=X\cup Y\) and \(X\cap Y=\alpha.\beta.\gamma\).
The set \(X\cup Y\) is the disjoint union of \(X\) and \(Y\setminus X\), and
\(Y\) is the disjoint union of \(Y\setminus X\) and \(X\cap Y\). Hence
\(N(X\cup Y)=N(X)+N(Y)-N(X\cap Y)\). Dividing by \(n\) gives the general
addition theorem, and \(N(\alpha.\beta.\gamma)=0\) removes the last term.

Item~5. Divide the first form of item~2 by
\({}_{\alpha}F''(\beta)>0\).

Item~6. Since \(N(\alpha.\gamma)>0\), item~5 with the roles of \(\beta\)
and \(\gamma\) exchanged gives
\({}_{\alpha.\gamma}F''(\beta_i)={}_{\alpha}F''(\beta_i.\gamma)/{}_{\alpha}F''(\gamma)\).
Since \(N(\alpha.\beta_i)>0\), item~2 gives
\({}_{\alpha}F''(\beta_i.\gamma)={}_{\alpha}F''(\beta_i)\cdot{}_{\alpha.\beta_i}F''(\gamma)\)
for the numerator. For the denominator, \(\alpha.\gamma\subset\Sigma\beta_i\)
means every element of \(\alpha.\gamma\) belongs to some \(\beta_j\), so
\(\alpha.\gamma\) is the union of the classes \(\alpha.\beta_j.\gamma\).
For \(i\neq j\) the classes \(\alpha.\beta_i.\gamma\) and
\(\alpha.\beta_j.\gamma\) are disjoint, because their intersection lies in
the empty class \(\alpha.\beta_i.\beta_j\). Hence
\(N(\alpha.\gamma)=\sum_jN(\alpha.\beta_j.\gamma)\), so
\({}_{\alpha}F''(\gamma)=\sum_j{}_{\alpha}F''(\beta_j.\gamma)
=\sum_j{}_{\alpha.\beta_j}F''(\gamma)\cdot{}_{\alpha}F''(\beta_j)\),
by item~2 for each \(j\). Substituting both expressions gives the formula.

Item~7. If \(\gamma\subset\beta\), then \(\alpha.\beta.\gamma=\alpha.\gamma\)
and \(\alpha.\beta\) is the disjoint union of \(\alpha.\gamma\) and
\(\alpha.\beta.\bar\gamma\). Hence
\(N(\alpha.\beta)=N(\alpha.\gamma)+N(\alpha.\beta.\bar\gamma)\geq N(\alpha.\gamma)\),
with equality exactly when \(N(\alpha.\beta.\bar\gamma)=0\). Divide by \(n\).
\end{proof}
\end{annexitem}

\begin{annexitem}{annex:p15-w02b-p15-027-sym}{thm:p15-w02b-p15-027-triple-symmetry}{Five buttons}
\begin{proof}
The description fixes the buttons up to naming. Call the round ones
\(b_1\) (black, not large), \(b_2\) (black, large), \(b_3\) (not black,
large) and \(b_4\) (neither). These are the four round buttons, with two
black, two large and one both. Call the square one \(b_5\) (black,
large). Then
\(\alpha=\{b_1,b_2,b_3,b_4\}\), \(\beta=\{b_1,b_2,b_5\}\) and
\(\gamma=\{b_2,b_3,b_5\}\). So \(\alpha.\beta=\{b_1,b_2\}\),
\(\alpha.\gamma=\{b_2,b_3\}\), \(\alpha.\beta.\gamma=\{b_2\}\) and
\(\beta.\gamma=\{b_2,b_5\}\).

Within \(\alpha\):
\({}_{\alpha}F''(\gamma)=2/4=1/2\) and
\({}_{\alpha.\beta}F''(\gamma)=N(\{b_2\})/N(\{b_1,b_2\})=1/2\). So
\(\gamma\) is independent of \(\beta\) within \(\alpha\). Because
\(N(\alpha.\beta)\) and \(N(\alpha.\gamma)\) are positive, item~3 of
\cref{thm:p15-w02b-p15-027-finite-calculus} gives mutual independence.

Within \(\beta\): \(\beta.\alpha=\alpha.\beta\), so
\({}_{\beta.\alpha}F''(\gamma)={}_{\alpha.\beta}F''(\gamma)=1/2\), while
\({}_{\beta}F''(\gamma)=N(\beta.\gamma)/N(\beta)=2/3\). The two differ, so
\(\gamma\) is not insensitive within \(\beta\) to selection of
\(\alpha\)-elements.

Finally, if \(\beta\subset\alpha\) then \(\beta.\alpha=\beta\), and the two
sides of \({}_{\beta.\alpha}F''(\gamma)={}_{\beta}F''(\gamma)\) are the same
quotient \(N(\beta.\gamma)/N(\beta)\).
\end{proof}
\end{annexitem}

\begin{annexitem}{annex:p15-w02b-p15-027-binomial}{thm:p15-w02b-p15-027-binomial}{Cyclic counts, Schiff's reduction and the binomial formula}
\begin{proof}
Two counting identities are used. Let \(\tau\) be a word of length
\(k\geq1\) and \(r\geq1\).
\begin{enumerate}
\item[(L)] \(c_a(\tau)=\sum_{\pi}c_a(\pi\tau)\), the sum over all words
  \(\pi\) of length \(r\).
\item[(R)] \(c_a(\tau)=\sum_{\pi}c_a(\tau\pi)\), the sum over all words
  \(\pi\) of length \(r\).
\end{enumerate}
For (L), an occurrence of \(\tau\) starting at position \(i\) determines
the word \(\pi=a_{i-r}\cdots a_{i-1}\) (indices modulo \(N\)), and then
\(\pi\tau\) occurs at \(i-r\). Conversely an occurrence of \(\pi\tau\) at
\(j\) gives an occurrence of \(\tau\) at \(j+r\). These maps are mutually
inverse, so occurrences of \(\tau\) correspond one to one with pairs of a
word \(\pi\) and an occurrence of \(\pi\tau\). (R) is proved in the same
way with the \(r\) terms after the occurrence. In particular (R) with
\(r=1\) gives \(c_a(\tau)=c_a(\tau0)+c_a(\tau1)\), so \(c_a(\tau b)\leq
c_a(\tau)\) for \(b\in\{0,1\}\). Also \(c_a(1)+c_a(0)=N\), so \(p+q=1\).

Item~1. The hypothesis is the defining condition for words of length
\(n\). Let \(\tau\) have length \(j<n\) and \(c_a(\tau)>0\). By (L) with
\(r=n-j\),
\(c_a(\tau)=\sum_\pi c_a(\pi\tau)\) and
\(c_a(\tau1)=\sum_\pi c_a(\pi\tau1)\), both sums over words \(\pi\) of
length \(n-j\). Each \(\pi\tau\) has length \(n\). If
\(c_a(\pi\tau)>0\), the hypothesis gives
\(c_a(\pi\tau1)=p\,c_a(\pi\tau)\). If \(c_a(\pi\tau)=0\), then
\(c_a(\pi\tau1)\leq c_a(\pi\tau)=0\), so again
\(c_a(\pi\tau1)=p\,c_a(\pi\tau)\). Summing over \(\pi\) gives
\(c_a(\tau1)=p\,c_a(\tau)\). Hence \(a\) is cyclically \(n\)-free.

Item~2. We show by induction on \(k\), for \(1\leq k\leq n\), that every
word \(w\) of length \(k\) with \(m\) ones has \(c_a(w)=Np^mq^{k-m}\).
For \(k=1\), \(c_a(1)=Np\) and \(c_a(0)=Nq\) by the definition of \(p,q\).
Let \(1\leq k\leq n-1\), assume the claim for length \(k\), and let
\(w=\tau b\) with \(\tau\) of length \(k\) having \(m'\) ones. Since
\(k\leq n-1\) and \(a\) is cyclically \((n-1)\)-free,
\(c_a(\tau)>0\) gives \(c_a(\tau1)=p\,c_a(\tau)\) and then, by (R),
\(c_a(\tau0)=c_a(\tau)-c_a(\tau1)=q\,c_a(\tau)\). If \(c_a(\tau)=0\), then
\(c_a(\tau b)=0=c_a(\tau)\cdot p\) and \(=c_a(\tau)\cdot q\). In either
case \(c_a(\tau1)=Np^{m'+1}q^{k-m'}\) and \(c_a(\tau0)=Np^{m'}q^{k+1-m'}\),
which is the claim for length \(k+1\). (For \(n=1\) only the base case is
needed.) Taking \(k=n\) gives
\({}_{a_{(n)}}F''(\sigma_m)=c_a(\sigma_m)/N=p^mq^{n-m}\). Each of the \(N\)
cyclic windows of length \(n\) is exactly one word, and the words of
length \(n\) with \(m\) ones are \(\binom{n}{m}\) in number, each with
\(c_a\)-value \(Np^mq^{n-m}\). So the number of windows with exactly
\(m\) ones is \(\binom{n}{m}Np^mq^{n-m}\), and dividing by \(N\) gives
the binomial formula.

Item~3. Cyclic \(n\)-freedom includes cyclic \(((n+1)-1)\)-freedom, so
item~2 with \(n+1\) in place of \(n\) and \(p=q=1/2\) gives
\(c_a(w)=N/2^{n+1}\) for every word \(w\) of length \(n+1\). This is a
nonnegative integer and positive because \(N\geq1\), so it is at least 1
and \(2^{n+1}\) divides \(N\); in particular \(N\geq2^{n+1}>n\). Let
\(b=a_0\cdots a_{N-1}a_0\cdots a_{n-1}\), so \(b_t=a_{t\bmod N}\) for
\(0\leq t\leq N+n-1\). The cyclic window of length \(n+1\) starting at
\(i\in\{0,\ldots,N-1\}\) consists of \(a_{(i+s)\bmod N}=b_{i+s}\) for
\(s=0,\ldots,n\), since \(i+s\leq N+n-1\). So every cyclic occurrence is
a block of \(b\), and every word of length \(n+1\), occurring cyclically
at least once, is a block of \(b\).

Item~4. Let \(N=2^{n+1}\) and each word of length \(n+1\) occur exactly
once. For a word \(w\) of length \(k\leq n+1\), (R) with \(r=n+1-k\)
(nothing to do when \(k=n+1\)) gives \(c_a(w)=2^{n+1-k}\), the number
of words \(\pi\) of length \(n+1-k\). Hence \(c_a(1)=2^n=N/2\), so
\(p=q=1/2\). For \(\tau\) of length \(j\leq n\),
\(c_a(\tau1)=2^{n-j}=\tfrac12\cdot2^{n+1-j}=p\,c_a(\tau)\), so \(a\) is
cyclically \(n\)-free.

The four periods. Their cyclic windows, starting at the first term, are
as follows. For (A), with length 2: 11, 10, 00, 01. For (B), with
length 3: 101, 011, 111, 110, 100, 000, 001, 010. For (C), with
length 4: 1011, 0110, 1100, 1000, 0000, 0001, 0011, 0111, 1111, 1110,
1101, 1010, 0100, 1001, 0010, 0101. For (D), with length 5: 01100,
11000, 10001, 00011, 00111, 01110, 11101, 11010, 10101, 01010, 10100,
01001, 10010, 00100, 01000, 10000, 00000, 00001, 00010, 00101, 01011,
10111, 01111, 11111, 11110, 11100, 11001, 10011, 00110, 01101, 11011,
10110.
In each list the words are pairwise distinct and their number is
\(2^{n+1}\) for \(n=1,2,3,4\) respectively, which is the number of all
words of that length. So each occurs exactly once, and item~4 gives
cyclic 1-, 2-, 3- and 4-freedom with \(p=1/2\). For (A) and the pair
\(\tau=11\): \(c_a(11)=1\), and its only occurrence, starting at the
first term, is followed by \(0\), so \(c_a(111)=0\neq\tfrac12\cdot1\);
(A) is not cyclically 2-free. For (B), item~2 with \(n=3\) gives
\(c_a(w)=8\cdot2^{-3}=1\) for each word of length 3, so
\({}_{a_{(3)}}F''(m)=\binom3m/8\).
\end{proof}
\end{annexitem}

\begin{annexitem}{annex:p15-w02b-p15-027-bernoulli-bound}{thm:p15-w02b-p15-027-bernoulli-bound}{Chebyshev's bound for the binomial weights}
\begin{proof}
Write \(B(m)=\binom nm p^mq^{n-m}\) for \(0\leq m\leq n\). By the binomial
theorem \(\sum_mB(m)=(p+q)^n=1\).

First moment. For \(m\geq1\),
\(m\binom nm=n\binom{n-1}{m-1}\). Hence
\[
 \sum_{m=0}^nmB(m)=np\sum_{m=1}^{n}\binom{n-1}{m-1}p^{m-1}q^{n-m}
 =np(p+q)^{n-1}=np.
\]
Second factorial moment. For \(n\geq2\) and \(m\geq2\),
\(m(m-1)\binom nm=n(n-1)\binom{n-2}{m-2}\), so
\(\sum_mm(m-1)B(m)=n(n-1)p^2(p+q)^{n-2}=n(n-1)p^2\). For \(n=1\) every
term has \(m(m-1)=0\) and the value \(n(n-1)p^2=0\) agrees.
Hence
\[
\begin{aligned}
 \sum_m(m-np)^2B(m)
 &=\sum_m m(m-1)B(m)+\sum_m mB(m)\\
 &\qquad-2np\sum_m mB(m)+n^2p^2\sum_mB(m)\\
 &=n(n-1)p^2+np-2n^2p^2+n^2p^2=npq.
\end{aligned}
\]
If \(|m/n-p|\geq\delta\) then \((m-np)^2\geq n^2\delta^2\). All terms
being nonnegative,
\[
 n^2\delta^2\sum_{|m/n-p|\geq\delta}B(m)\leq\sum_m(m-np)^2B(m)=npq ,
\]
which is the displayed inequality after division by \(n^2\delta^2\).

Item~1. By item~2 of \cref{thm:p15-w02b-p15-027-binomial}, the
proportion of cyclic \(n\)-segments with exactly \(m\) ones is \(B(m)\).
A segment is not \(\Delta p\) exactly when its \(m\) satisfies
\(|m/n-p|\geq\delta\). Summing over those \(m\) gives a proportion
\(\sum_{|m/n-p|\geq\delta}B(m)\leq pq/(n\delta^2)\).

Item~2. The properties \(m=0,\ldots,n\) of the segments are mutually
exclusive and exhaust them, and \(\Delta p\) is the union of those
\(m\) with \(|m/n-p|<\delta\). On a finite initial part of the segment
sequence, the relative frequency of \(\Delta p\) is the sum of the
relative frequencies of those \(m\). Formula (3) asserts that each
summand converges to \(B(m)\), so the sum converges to
\(\sum_{|m/n-p|<\delta}B(m)=1-\sum_{|m/n-p|\geq\delta}B(m)\geq1-pq/(n\delta^2)\).
A relative frequency is at most 1, so
\(1-pq/(n\delta^2)\leq{}_{\alpha_n}F(\Delta p)\leq1\) for every \(n\).
Since \(pq/(n\delta^2)\to0\) as \(n\to\infty\), the limit is 1.

Item~3. With \(0<p<1\) we have \(pq>0\), so
\(\delta_\varepsilon=\sqrt{pq/(n\varepsilon)}>0\) and the main inequality
with \(\delta=\delta_\varepsilon\) bounds the sum by
\(pq/(n\delta_\varepsilon^2)=\varepsilon\). If
\(n\geq pq/(\varepsilon_0\varphi^2)\) and \(\varepsilon\geq\varepsilon_0\),
then \(\delta_\varepsilon^2=pq/(n\varepsilon)\leq pq/(n\varepsilon_0)\leq\varphi^2\),
so \(\delta_\varepsilon\leq\varphi\).
\end{proof}
\end{annexitem}

\begin{annexitem}{annex:p15-w02b-p15-027-middle-frequencies}{thm:p15-w02b-p15-027-middle-frequencies}{Accumulation points of a frequency sequence}
\begin{proof}
Write \(S_n=a_1+\cdots+a_n\), so \(h_n=S_n/n\in[0,1]\). A number \(c\)
is an accumulation point (middle frequency) when for every
\(\varepsilon>0\) and every \(K\) there is \(n>K\) with
\(|h_n-c|<\varepsilon\). Two facts are used.
First, \(h_{n+1}-h_n=(nS_n+na_{n+1}-(n+1)S_n)/(n(n+1))=(a_{n+1}-h_n)/(n+1)\),
and \(|a_{n+1}-h_n|\leq1\), so \(|h_{n+1}-h_n|\leq1/(n+1)\).
Second, let \(L=\liminf_nh_n\) and \(U=\limsup_nh_n\), where
\(U=\lim_{K\to\infty}u_K\) with \(u_K=\sup_{n>K}h_n\) nonincreasing, and
\(L\) is defined dually.

Item~1. The sequence \((h_n)\) is bounded, so by the Bolzano--Weierstrass
theorem a subsequence converges to some \(c\). Every neighbourhood of
\(c\) then contains \(h_n\) for infinitely many \(n\), so \(c\) is an
accumulation point. If \(c\notin[0,1]\), then the positive distance from
\(c\) to \([0,1]\) bounds \(|h_n-c|\) from below, so no point outside
\([0,1]\) is an accumulation point.

Item~2. \(U\) is an accumulation point: given \(\varepsilon>0\) and \(K\),
choose \(K'\geq K\) with \(u_{K'}<U+\varepsilon\); since
\(u_{K'}\geq U\), there is \(n>K'\) with \(h_n>u_{K'}-\varepsilon\geq U-\varepsilon\),
and \(h_n\leq u_{K'}<U+\varepsilon\). The same argument for lower bounds
shows that \(L\) is an accumulation point. No accumulation point exceeds
\(U\): if \(c>U\), put \(\varepsilon=(c-U)/2\) and choose \(K\) with
\(u_K<U+\varepsilon\); then \(h_n<c-\varepsilon\) for every \(n>K\).
Dually none is below \(L\). So the set \(M\) lies in \([L,U]\) and
contains both endpoints.
Now let \(L<c<U\), \(\varepsilon>0\) and \(K\) be given. Choose
\(K_1\geq K\) with \(1/(K_1+1)<\varepsilon\). As \(U\) is an
accumulation point and \(U-c>0\), there is \(n_1>K_1\) with
\(h_{n_1}>c\). As \(L\) is an accumulation point and \(c-L>0\), there is
\(n_2>n_1\) with \(h_{n_2}<c\). Let \(n^*\) be the least \(n\) with
\(n_1<n\leq n_2\) and \(h_n\leq c\); it exists since \(n_2\) qualifies.
Then \(h_{n^*-1}>c\geq h_{n^*}\), whence
\(0\leq c-h_{n^*}<h_{n^*-1}-h_{n^*}\leq1/n^*\leq1/(K_1+1)<\varepsilon\),
with \(n^*>K\). So \(c\in M\), and \(M=[L,U]\). If \(M\) has more than one
element then \(L<U\) and \(M\) is a nondegenerate interval.

Item~3. If \(h_n\to l\), every convergent subsequence has limit \(l\),
and an accumulation point is the limit of a subsequence chosen with
\(|h_{n_j}-c|<1/j\) and increasing \(n_j\); so \(M=\{l\}\). Conversely,
if \(M=\{l\}\), item~2 gives \(L=U=l\). Given \(\varepsilon>0\) choose
\(K\) with \(u_K<l+\varepsilon\) and, dually,
\(\inf_{n>K}h_n>l-\varepsilon\); then \(|h_n-l|<\varepsilon\) for
\(n>K\), so \(h_n\to l\), and \(l\) is the frequency limit.
\end{proof}
\end{annexitem}

\begin{annexitem}{annex:p15-w02b-p15-027-finite-records}{thm:p15-w02b-p15-027-finite-records}{Altering a finite initial record; bounded and unbounded existentials}
\begin{proof}
Item~1. Let \(a\) be absolutely free with probability \(p\) and let \(w\)
have length \(\ell\). Define \(a'\) by \(a'_t=w_t\) for \(t\leq\ell\) and
\(a'_t=a_t\) for \(t>\ell\). The numbers of ones among the first \(n\)
terms of \(a\) and \(a'\) differ by at most \(\ell\), so their relative
frequencies differ by at most \(\ell/n\), and \(a'\) has the frequency
limit \(p\).

Fix a word \(\tau\) of length \(k\). A position \(t>\ell+k\) has
\(t-k>\ell\), so \(a'\) and \(a\) agree at \(t-k,\ldots,t\). Hence
\(t\in S'_\tau\) iff \(t\in S_\tau\) (primes for \(a'\)), and the term at
\(t\) is the same. Let \(s(n),s'(n)\) be the numbers of selected
positions up to \(n\), and \(o(n),o'(n)\) the numbers of selected ones.
Then \(|s(n)-s'(n)|\leq\ell+k\) and \(|o(n)-o'(n)|\leq\ell+k\). If
\(S'_\tau\) is infinite, so is \(S_\tau\), since the two differ only at
positions up to \(\ell+k\). For \(n\) with \(s(n),s'(n)>0\), using
\(o(n)\leq s(n)\),
\[
 \Bigl|\frac{o'(n)}{s'(n)}-\frac{o(n)}{s(n)}\Bigr|
 \leq\frac{|o'(n)-o(n)|}{s'(n)}+o(n)\frac{|s(n)-s'(n)|}{s(n)s'(n)}
 \leq\frac{2(\ell+k)}{s'(n)} .
\]
Taking \(n\) as the position of the \(r\)th element of \(S'_\tau\), the
left fraction \(o'(n)/s'(n)\) is the relative frequency of ones among the
first \(r\) selected elements of \(a'\). As \(r\to\infty\), \(s'(n)=r\to\infty\)
and \(s(n)\to\infty\). By absolute freedom of \(a\),
\(o(n)/s(n)\to p\) along these \(n\), because these are relative
frequencies of the selected subsequence of \(a\) at increasing lengths.
The bound then gives \(o'(n)/s'(n)\to p\). So \(a'\) is absolutely free
with probability \(p\) and begins with \(w\).

Item~2. Let \(w\) have length \(\ell\). The alternative \(w\) followed
by zeros satisfies \(S_\infty\): for any \(x\) take
\(y=\max(x,\ell)+1\), so \(y>x\) and \(a_y=0\). The alternative \(w\)
followed by ones violates \(S_\infty\): for \(x=\ell\) every \(y>x\)
has \(a_y=1\). If \(a_1=a_2=1\), then \(x=1\) gives \(a_x=1\) and
\(a_{x+1}\neq0\), so \(S_b\) fails.
\end{proof}
\end{annexitem}

\begin{annexitem}{annex:p15-w02b-p15-027-refs}{thm:p15-w02b-p15-027-reference-class}{Reference classes}
\begin{proof}
Item~1. Since \(\alpha.\beta.\beta=\alpha.\beta\) and
\(N(\alpha.\beta)>0\), \({}_{\alpha.\beta}F''(\beta)=1\). By item~1 of
\cref{thm:p15-w02b-p15-027-finite-calculus},
\(N(\alpha.\bar\beta)=N(\alpha)-N(\alpha.\beta)>0\), and
\(\alpha.\bar\beta.\beta\) is empty, so
\({}_{\alpha.\bar\beta}F''(\beta)=0\). The value
\({}_{\alpha}F''(\beta)=N(\alpha.\beta)/N(\alpha)\) lies strictly between
0 and 1 by the hypotheses, and its definition mentions no particular
element of \(\alpha\). Both \(\alpha.\beta\) and \(\alpha.\bar\beta\) are
nonempty.

Item~2. The element \(1\) belongs to \(\alpha.\beta\) and not to
\(\gamma\), so the universal law fails. Among the first \(n\) positive
integers all belong to \(\beta\) and \(n-1\) to \(\gamma\); the relative
frequencies \(1\) and \((n-1)/n\) both converge to \(1\).
\end{proof}
\end{annexitem}
